% Options for packages loaded elsewhere
% Options for packages loaded elsewhere
\PassOptionsToPackage{unicode}{hyperref}
\PassOptionsToPackage{hyphens}{url}
%
\documentclass[
  ignorenonframetext,
  aspectratio=169,
  russian,
  aspectratio=169]{beamer}
\newif\ifbibliography
\usepackage{pgfpages}
\setbeamertemplate{caption}[numbered]
\setbeamertemplate{caption label separator}{: }
\setbeamercolor{caption name}{fg=normal text.fg}
\beamertemplatenavigationsymbolshorizontal
% remove section numbering
\setbeamertemplate{part page}{
  \centering
  \begin{beamercolorbox}[sep=16pt,center]{part title}
    \usebeamerfont{part title}\insertpart\par
  \end{beamercolorbox}
}
\setbeamertemplate{section page}{
  \centering
  \begin{beamercolorbox}[sep=12pt,center]{section title}
    \usebeamerfont{section title}\insertsection\par
  \end{beamercolorbox}
}
\setbeamertemplate{subsection page}{
  \centering
  \begin{beamercolorbox}[sep=8pt,center]{subsection title}
    \usebeamerfont{subsection title}\insertsubsection\par
  \end{beamercolorbox}
}
% Prevent slide breaks in the middle of a paragraph
\widowpenalties 1 10000
\raggedbottom
\AtBeginPart{
  \frame{\partpage}
}
\AtBeginSection{
  \ifbibliography
  \else
    \frame{\sectionpage}
  \fi
}
\AtBeginSubsection{
  \frame{\subsectionpage}
}
\usepackage{iftex}
\ifPDFTeX
  \usepackage[T1]{fontenc}
  \usepackage[utf8]{inputenc}
  \usepackage{textcomp} % provide euro and other symbols
\else % if luatex or xetex
  \usepackage{unicode-math} % this also loads fontspec
  \defaultfontfeatures{Scale=MatchLowercase}
  \defaultfontfeatures[\rmfamily]{Ligatures=TeX,Scale=1}
\fi
\usepackage{lmodern}

\usefonttheme{serif} % use mainfont rather than sansfont for slide text
\ifPDFTeX\else
  % xetex/luatex font selection
  \setmainfont[]{DejaVu Serif}
  \setsansfont[]{DejaVu Sans}
  \setmonofont[]{DejaVu Sans Mono}
\fi
% Use upquote if available, for straight quotes in verbatim environments
\IfFileExists{upquote.sty}{\usepackage{upquote}}{}
\IfFileExists{microtype.sty}{% use microtype if available
  \usepackage[]{microtype}
  \UseMicrotypeSet[protrusion]{basicmath} % disable protrusion for tt fonts
}{}


\usepackage{longtable,booktabs,array}
\usepackage{calc} % for calculating minipage widths
\usepackage{caption}
% Make caption package work with longtable
\makeatletter
\def\fnum@table{\tablename~\thetable}
\makeatother
\usepackage{graphicx}
\makeatletter
\newsavebox\pandoc@box
\newcommand*\pandocbounded[1]{% scales image to fit in text height/width
  \sbox\pandoc@box{#1}%
  \Gscale@div\@tempa{\textheight}{\dimexpr\ht\pandoc@box+\dp\pandoc@box\relax}%
  \Gscale@div\@tempb{\linewidth}{\wd\pandoc@box}%
  \ifdim\@tempb\p@<\@tempa\p@\let\@tempa\@tempb\fi% select the smaller of both
  \ifdim\@tempa\p@<\p@\scalebox{\@tempa}{\usebox\pandoc@box}%
  \else\usebox{\pandoc@box}%
  \fi%
}
% Set default figure placement to htbp
\def\fps@figure{htbp}
\makeatother



\ifLuaTeX
\usepackage[bidi=basic,provide=*]{babel}
\else
\usepackage[bidi=default,provide=*]{babel}
\fi
\ifPDFTeX
\else
\babelfont{rm}[]{DejaVu Serif}
\fi
% get rid of language-specific shorthands (see #6817):
\let\LanguageShortHands\languageshorthands
\def\languageshorthands#1{}


\setlength{\emergencystretch}{3em} % prevent overfull lines

\providecommand{\tightlist}{%
  \setlength{\itemsep}{0pt}\setlength{\parskip}{0pt}}



 

\usepackage[]{csquotes}

\IfFileExists{plex-otf.sty}{
  %% Full TeXlive
  % \usepackage[%
  %   % math,
  %   RM={Scale=0.94},SS={Scale=0.94},SScon={Scale=0.94},TT={Scale=MatchLowercase,FakeStretch=0.9},DefaultFeatures={Ligatures=Common}
  % ]{plex-otf}
}{
  %% TinyTeX
  \usepackage{libertine}
}

%%% Load theme
% https://deic.uab.cat/~iblanes/beamer_gallery/
\IfFileExists{beamerthemegotham.sty}{
  %% Full TeXlive
  \usetheme{gotham}
  \gothamset{
    numbering=totalpagenumber,
    parttocframe default=off,
    sectiontocframe default=off,
    subsectiontocframe default=off,
  }
}{
  %% TinyTeX
  \usetheme{Madrid}
}
\makeatletter
\@ifpackageloaded{caption}{}{\usepackage{caption}}
\AtBeginDocument{%
\ifdefined\contentsname
  \renewcommand*\contentsname{Содержание}
\else
  \newcommand\contentsname{Содержание}
\fi
\ifdefined\listfigurename
  \renewcommand*\listfigurename{Список иллюстраций}
\else
  \newcommand\listfigurename{Список иллюстраций}
\fi
\ifdefined\listtablename
  \renewcommand*\listtablename{Список таблиц}
\else
  \newcommand\listtablename{Список таблиц}
\fi
\ifdefined\figurename
  \renewcommand*\figurename{Рисунок}
\else
  \newcommand\figurename{Рисунок}
\fi
\ifdefined\tablename
  \renewcommand*\tablename{Таблица}
\else
  \newcommand\tablename{Таблица}
\fi
}
\@ifpackageloaded{float}{}{\usepackage{float}}
\floatstyle{ruled}
\@ifundefined{c@chapter}{\newfloat{codelisting}{h}{lop}}{\newfloat{codelisting}{h}{lop}[chapter]}
\floatname{codelisting}{Список}
\newcommand*\listoflistings{\listof{codelisting}{Листинги}}
\makeatother
\makeatletter
\makeatother
\makeatletter
\@ifpackageloaded{caption}{}{\usepackage{caption}}
\@ifpackageloaded{subcaption}{}{\usepackage{subcaption}}
\makeatother

\usepackage{bookmark}
\IfFileExists{xurl.sty}{\usepackage{xurl}}{} % add URL line breaks if available
\urlstyle{same}
\hypersetup{
  pdftitle={Выполнение лабораторной работы №4},
  pdfauthor={Карпухин Клим},
  pdflang={ru-RU},
  hidelinks,
  pdfcreator={LaTeX via pandoc}}


\title{Выполнение лабораторной работы №4}
\subtitle{Продвинутое использование git}
\author{Карпухин Клим}
\date{2026-03-07}

\begin{document}
\frame{\titlepage}


\section{Содержание}\label{ux441ux43eux434ux435ux440ux436ux430ux43dux438ux435}

\begin{frame}{Содержание}
\begin{enumerate}[<+->]
\tightlist
\item
  Информация о докладчике
\item
  Вводная часть и актуальность
\item
  Объект и предмет исследования
\item
  Научная новизна и практическая значимость
\item
  Цель, гипотеза и задачи
\item
  Материалы, методы и инструменты
\item
  Ход работы (этапы, скриншоты)
\item
  Результаты и анализ
\item
  Выводы
\end{enumerate}
\end{frame}

\section{Информация}\label{ux438ux43dux444ux43eux440ux43cux430ux446ux438ux44f}

\begin{frame}{Докладчик}
\phantomsection\label{ux434ux43eux43aux43bux430ux434ux447ux438ux43a}
\begin{columns}[c]
\begin{column}{0.65\linewidth}
\begin{itemize}[<+->]
\tightlist
\item
  \textbf{Карпухин Клим}
\item
  Российский университет дружбы народов
\item
  Email:
  \href{mailto:1032255580@rudn.ru}{\nolinkurl{1032255580@rudn.ru}}
\item
  Роли: студент (лабораторная работа по ОС/виртуализации)
\end{itemize}
\end{column}

\begin{column}{0.35\linewidth}

\includegraphics[width=0.9\linewidth,height=\textheight,keepaspectratio]{./image/dokladchik.png}
\end{column}
\end{columns}
\end{frame}

\section{Вводная
часть}\label{ux432ux432ux43eux434ux43dux430ux44f-ux447ux430ux441ux442ux44c}

\begin{frame}{Актуальность}
\phantomsection\label{ux430ux43aux442ux443ux430ux43bux44cux43dux43eux441ux442ux44c}
\begin{itemize}[<+->]
\tightlist
\item
  Современная разработка программного обеспечения требует чёткой
  организации веток, релизов и истории изменений. Gitflow решает эту
  задачу.
\item
  Семантическое версионирование и общепринятые коммиты облегчают
  автоматизацию выпуска версий и генерации журналов изменений.
\item
  Владение такими инструментами делает процесс разработки прозрачным и
  воспроизводимым, что важно и в учебных, и в реальных проектах.
\end{itemize}
\end{frame}

\begin{frame}{Объект и предмет исследования}
\phantomsection\label{ux43eux431ux44aux435ux43aux442-ux438-ux43fux440ux435ux434ux43cux435ux442-ux438ux441ux441ux43bux435ux434ux43eux432ux430ux43dux438ux44f}
\begin{itemize}[<+->]
\tightlist
\item
  \textbf{Объект:} процесс управления репозиторием программного проекта
  с использованием системы контроля версий git.
\item
  \textbf{Предмет:} методология Gitflow, семантическое версионирование
  (SemVer) и спецификация Conventional Commits, а также инструменты
  commitizen и standard-changelog.
\end{itemize}
\end{frame}

\section{Научная новизна и практическая
значимость}\label{ux43dux430ux443ux447ux43dux430ux44f-ux43dux43eux432ux438ux437ux43dux430-ux438-ux43fux440ux430ux43aux442ux438ux447ux435ux441ux43aux430ux44f-ux437ux43dux430ux447ux438ux43cux43eux441ux442ux44c}

\begin{frame}{Научная новизна}
\phantomsection\label{ux43dux430ux443ux447ux43dux430ux44f-ux43dux43eux432ux438ux437ux43dux430}
\begin{itemize}[<+->]
\tightlist
\item
  Интеграция Gitflow, SemVer и Conventional Commits в единый
  практический сценарий на учебном репозитории.
\item
  Использование стандартизированных сообщений коммитов как источника
  данных для автоматической генерации CHANGELOG.
\item
  Формирование практических навыков, необходимых для внедрения
  промышленного процесса релизов даже в небольших учебных проектах.
\end{itemize}
\end{frame}

\begin{frame}{Практическая значимость работы}
\phantomsection\label{ux43fux440ux430ux43aux442ux438ux447ux435ux441ux43aux430ux44f-ux437ux43dux430ux447ux438ux43cux43eux441ux442ux44c-ux440ux430ux431ux43eux442ux44b}
\begin{itemize}[<+->]
\tightlist
\item
  Освоенные подходы позволяют эффективно управлять ветками feature,
  release и hotfix в реальных командах разработки.
\item
  Автоматическая генерация журналов изменений и релизов снижает
  количество ручной работы и ошибок при подготовке версий.
\item
  Полученные навыки можно напрямую перенести в курсовые, дипломные и
  практические проекты с использованием GitHub.
\end{itemize}
\end{frame}

\section{Цель, гипотеза и
задачи}\label{ux446ux435ux43bux44c-ux433ux438ux43fux43eux442ux435ux437ux430-ux438-ux437ux430ux434ux430ux447ux438}

\begin{frame}{Цель}
\phantomsection\label{ux446ux435ux43bux44c}
Получить устойчивые навыки корректной работы с репозиториями git на
основе методологии Gitflow и практик семантического версионирования и
общепринятых коммитов.
\end{frame}

\begin{frame}{Гипотеза}
\phantomsection\label{ux433ux438ux43fux43eux442ux435ux437ux430}
Если применить Gitflow, SemVer и Conventional Commits в учебном
репозитории, то процесс разработки и выпуска версий станет более
структурированным и прозрачным.
\end{frame}

\begin{frame}{Задачи}
\phantomsection\label{ux437ux430ux434ux430ux447ux438}
\begin{itemize}[<+->]
\tightlist
\item
  Установить и настроить инструменты git-flow, Node.js, pnpm, commitizen
  и standard-changelog.
\item
  Создать и настроить репозиторий git-extended на GitHub.
\item
  Преобразовать репозиторий в репозиторий git-flow и настроить
  conventional commits.
\item
  Реализовать практический сценарий: разработка новой функциональности и
  выпуск релизов 1.0.0 и 1.2.3.
\end{itemize}
\end{frame}

\section{Материалы и
методы}\label{ux43cux430ux442ux435ux440ux438ux430ux43bux44b-ux438-ux43cux435ux442ux43eux434ux44b}

\begin{frame}{Материалы и методы}
\phantomsection\label{ux43cux430ux442ux435ux440ux438ux430ux43bux44b-ux438-ux43cux435ux442ux43eux434ux44b-1}
\begin{itemize}[<+->]
\tightlist
\item
  \textbf{Инструменты:}

  \begin{itemize}[<+->]
  \tightlist
  \item
    Система контроля версий \textbf{git}, расширение \textbf{git-flow}
    для поддержки Gitflow Workflow.
  \item
    Платформа \textbf{GitHub} для хранения и публикации репозитория
    git-extended.
  \item
    Среда \textbf{Node.js} и менеджер пакетов \textbf{pnpm} для
    установки commitizen и standard-changelog.
  \item
    Спецификация \textbf{Semantic Versioning (SemVer)} и
    \textbf{Conventional commits} как теоретическая база
    структурирования версий и коммитов.
  \end{itemize}
\end{itemize}
\end{frame}

\section{Ход работы --- подготовка и настройка
инструментов}\label{ux445ux43eux434-ux440ux430ux431ux43eux442ux44b-ux43fux43eux434ux433ux43eux442ux43eux432ux43aux430-ux438-ux43dux430ux441ux442ux440ux43eux439ux43aux430-ux438ux43dux441ux442ux440ux443ux43cux435ux43dux442ux43eux432}

\begin{frame}{Этап 1: Установка git-flow и настройка репозитория}
\phantomsection\label{ux44dux442ux430ux43f-1-ux443ux441ux442ux430ux43dux43eux432ux43aux430-git-flow-ux438-ux43dux430ux441ux442ux440ux43eux439ux43aux430-ux440ux435ux43fux43eux437ux438ux442ux43eux440ux438ux44f}
\begin{itemize}[<+->]
\tightlist
\item
  Установка git-flow из коллекции репозиториев copr.
\end{itemize}

\begin{figure}

\centering{

\includegraphics[width=0.3\linewidth,height=\textheight,keepaspectratio]{image/1.png}

}

\caption{\label{fig-001}Установка git-flow}

\end{figure}%
\end{frame}

\begin{frame}{Этап 2: Настройка Node.js и pnpm}
\phantomsection\label{ux44dux442ux430ux43f-2-ux43dux430ux441ux442ux440ux43eux439ux43aux430-node.js-ux438-pnpm}
\begin{itemize}[<+->]
\tightlist
\item
  Конфигурация Node.js и pnpm, добавление каталога с исполняемыми
  файлами в переменную PATH.
\end{itemize}

\begin{figure}

\centering{


\includegraphics[width=0.3\linewidth,height=\textheight,keepaspectratio]{image/2.png}

}

\caption{\label{fig-002}Добавление каталога pnpm в PATH}

\end{figure}%
\end{frame}

\begin{frame}{Этап 3: Установка commitizen и standard-changelog}
\phantomsection\label{ux44dux442ux430ux43f-3-ux443ux441ux442ux430ux43dux43eux432ux43aux430-commitizen-ux438-standard-changelog}
\begin{itemize}[<+->]
\tightlist
\item
  Установка глобальных пакетов для форматирования коммитов и генерации
  логов.
\end{itemize}

\includegraphics[width=0.3\linewidth,height=\textheight,keepaspectratio]{image/3.png}
\includegraphics[width=0.3\linewidth,height=\textheight,keepaspectratio]{image/4.png}
\end{frame}

\begin{frame}{Этап 4: Создание репозитория на GitHub}
\phantomsection\label{ux44dux442ux430ux43f-4-ux441ux43eux437ux434ux430ux43dux438ux435-ux440ux435ux43fux43eux437ux438ux442ux43eux440ux438ux44f-ux43dux430-github}
\begin{itemize}[<+->]
\tightlist
\item
  Создание репозитория git-extended на GitHub.
\end{itemize}

\begin{figure}

\centering{


\includegraphics[width=0.3\linewidth,height=\textheight,keepaspectratio]{image/5.png}

}

\caption{\label{fig-005}Создание репозитория git-extended}

\end{figure}%
\end{frame}

\begin{frame}{Этап 5: Первый коммит и отправка на GitHub}
\phantomsection\label{ux44dux442ux430ux43f-5-ux43fux435ux440ux432ux44bux439-ux43aux43eux43cux43cux438ux442-ux438-ux43eux442ux43fux440ux430ux432ux43aux430-ux43dux430-github}
\begin{itemize}[<+->]
\tightlist
\item
  Создание первого коммита и отправка его на GitHub.
\end{itemize}

\begin{figure}

\centering{

\includegraphics[width=0.3\linewidth,height=\textheight,keepaspectratio]{image/6.png}

}

\caption{\label{fig-006}Первый коммит и push}

\end{figure}%
\end{frame}

\section{Ход работы --- реализация сценария
Gitflow}\label{ux445ux43eux434-ux440ux430ux431ux43eux442ux44b-ux440ux435ux430ux43bux438ux437ux430ux446ux438ux44f-ux441ux446ux435ux43dux430ux440ux438ux44f-gitflow}

\begin{frame}{Этап 6: Инициализация git-flow}
\phantomsection\label{ux44dux442ux430ux43f-6-ux438ux43dux438ux446ux438ux430ux43bux438ux437ux430ux446ux438ux44f-git-flow}
\begin{itemize}[<+->]
\tightlist
\item
  Преобразование репозитория в репозиторий git-flow.
\end{itemize}
\end{frame}

\begin{frame}{Этап 7: Создание релиза 1.0.0}
\phantomsection\label{ux44dux442ux430ux43f-7-ux441ux43eux437ux434ux430ux43dux438ux435-ux440ux435ux43bux438ux437ux430-1.0.0}
\begin{itemize}[<+->]
\tightlist
\item
  Создание релиза с версией 1.0.0.
\end{itemize}

\begin{figure}

\centering{

\includegraphics[width=0.3\linewidth,height=\textheight,keepaspectratio]{image/16.png}

}

\caption{\label{fig-016}Создание релиза 1.0.0}

\end{figure}%
\end{frame}

\begin{frame}{Этап 8: Генерация журнала изменений}
\phantomsection\label{ux44dux442ux430ux43f-8-ux433ux435ux43dux435ux440ux430ux446ux438ux44f-ux436ux443ux440ux43dux430ux43bux430-ux438ux437ux43cux435ux43dux435ux43dux438ux439}
\begin{itemize}[<+->]
\tightlist
\item
  Создание журнала изменений с помощью standard-changelog.
\end{itemize}

\begin{figure}

\centering{

\includegraphics[width=0.3\linewidth,height=\textheight,keepaspectratio]{image/17.png}

}

\caption{\label{fig-017}Генерация CHANGELOG}

\end{figure}%
\end{frame}

\begin{frame}{Этап 9: Завершение релиза}
\phantomsection\label{ux44dux442ux430ux43f-9-ux437ux430ux432ux435ux440ux448ux435ux43dux438ux435-ux440ux435ux43bux438ux437ux430}
\begin{itemize}[<+->]
\tightlist
\item
  Отправка релизной ветки в основную ветку и публикация релиза на
  GitHub.
\end{itemize}

\begin{figure}

\centering{

\includegraphics[width=0.3\linewidth,height=\textheight,keepaspectratio]{image/19.png}

}

\caption{\label{fig-019}Завершение релиза}

\end{figure}%
\end{frame}

\begin{frame}{Этап 10: Создание релиза 1.2.3}
\phantomsection\label{ux44dux442ux430ux43f-10-ux441ux43eux437ux434ux430ux43dux438ux435-ux440ux435ux43bux438ux437ux430-1.2.3}
\begin{itemize}[<+->]
\tightlist
\item
  Повторение цикла для версии 1.2.3 (добавление новой функциональности).
\end{itemize}
\end{frame}

\section{Результаты и
анализ}\label{ux440ux435ux437ux443ux43bux44cux442ux430ux442ux44b-ux438-ux430ux43dux430ux43bux438ux437}

\begin{frame}{Анализ достигнутых результатов}
\phantomsection\label{ux430ux43dux430ux43bux438ux437-ux434ux43eux441ux442ux438ux433ux43dux443ux442ux44bux445-ux440ux435ux437ux443ux43bux44cux442ux430ux442ux43eux432}
\begin{itemize}[<+->]
\tightlist
\item
  Реализован полный цикл Gitflow: от создания функциональных веток до
  релизов и горячих исправлений.
\item
  Сообщения коммитов приведены к единому формату, что позволило
  автоматически форматировать CHANGELOG.
\item
  Репозиторий стал удобен для отслеживания истории и выпуска версий, что
  подтверждается успешно созданными релизами 1.0.0 и 1.2.3.
\end{itemize}
\end{frame}

\begin{frame}{Практическая значимость результатов}
\phantomsection\label{ux43fux440ux430ux43aux442ux438ux447ux435ux441ux43aux430ux44f-ux437ux43dux430ux447ux438ux43cux43eux441ux442ux44c-ux440ux435ux437ux443ux43bux44cux442ux430ux442ux43eux432}
\begin{itemize}[<+->]
\tightlist
\item
  Полученный опыт можно использовать при разработке учебных и реальных
  проектов с командной работой и релизным циклом.
\item
  Набор инструментов (git-flow, commitizen, standard-changelog)
  масштабируется на более крупные системы без существенных изменений
  процесса.
\item
  Сформированные навыки лягут в основу более сложных практик DevOps и
  CI/CD.
\end{itemize}
\end{frame}

\section{Выводы}\label{ux432ux44bux432ux43eux434ux44b}

\begin{frame}{Общее заключение}
\phantomsection\label{ux43eux431ux449ux435ux435-ux437ux430ux43aux43bux44eux447ux435ux43dux438ux435}
\begin{itemize}[<+->]
\tightlist
\item
  В ходе лабораторной работы освоены методология Gitflow, семантическое
  версионирование и стандарт общепринятых коммитов, а также инструменты
  commitizen и standard-changelog.
\item
  Практически отработаны сценарии создания функциональных веток,
  подготовки релизов и автоматической генерации журналов изменений.
\item
  Достигнутая цель подтверждает гипотезу о том, что применение данных
  практик делает процесс разработки более структурированным и
  прозрачным.
\end{itemize}
\end{frame}

\begin{frame}{Выводы}
\phantomsection\label{ux432ux44bux432ux43eux434ux44b-1}
\begin{enumerate}[<+->]
\tightlist
\item
  Gitflow обеспечивает удобную и понятную структуру веток для
  релизно-ориентированной разработки.
\item
  Семантическое версионирование и Conventional Commits упрощают
  сопровождение проекта и взаимодействие участников команды.
\item
  Автоматизация формирования CHANGELOG и релизов снижает вероятность
  ошибок и ускоряет выпуск новых версий.
\item
  Полученные компетенции полезны для последующей учебной,
  исследовательской и профессиональной деятельности в области разработки
  ПО.
\end{enumerate}
\end{frame}




\end{document}
